\documentclass[print,amsmath,amssymb,aps,prc,nofootinbib,notitlepage]{revtex4-1}

\usepackage[usenames,dvipsnames]{xcolor}
\usepackage{graphicx}% Include figure files
\usepackage{dcolumn}% Align table columns on decimal point
\usepackage{bm}% bold math
\usepackage{hyperref}% add hypertext capabilities
\usepackage[utf8]{inputenc}
\usepackage{microtype}

\renewcommand{\topfraction}{.85}
\renewcommand{\bottomfraction}{.7}
\renewcommand{\textfraction}{.15}
\renewcommand{\floatpagefraction}{.66}
\renewcommand{\dbltopfraction}{.66}
\renewcommand{\dblfloatpagefraction}{.66}
\setcounter{topnumber}{9}
\setcounter{bottomnumber}{9}
\setcounter{totalnumber}{20}
\setcounter{dbltopnumber}{9}

\usepackage{multirow}

\newcommand*{\w}[1]{\textcolor{red}{WR: #1}}
\newcommand*{\wrnote}[1]{\footnote{\textcolor{red}{WR: #1}}}
\newcommand{\phh}[1]{\textcolor{blue}{#1}}
\newcommand{\phhnote}[1]{\footnote{\textcolor{blue}{PHH: #1}}}
\newcommand{\nn}{\nonumber}
\renewcommand{\vec}[1]{{\bf #1}}
\newcommand{\rmd}{{\rm d}}
\newcommand{\MeV}{{\rm MeV}}
\newcommand{\fm}{{\rm fm}}
\newcommand{\iunit}{{\rm i}}
\newcommand{\econ}{{\rm e}}
\newcommand{\rvec}{\vec{r}}
\newcommand{\rvecp}{\vec{r}{\, }'}
\newcommand{\xvec}{\vec{x}}
\newcommand{\xvecp}{\vec{x}{\, }'}
\newcommand{\etal}{\emph{et al.}}

\newcommand{\tensym}[1]{\underline{\textsf{#1}}}
\newcommand{\nuc}[2]{$^{#1}$\textrm{#2}}

\renewcommand{\Re}{\text{Re}}
\renewcommand{\Im}{\text{Im}}

\usepackage{fancyvrb}

%%%%%%%%%%%%%%%%%%%%%%%%%%%%%%%%%%%%%%%%%%%%%%%%%%%%%%%%%%%%%%%%%%%%%%%%%%%%%%%%
\begin{document}

\author{W. Ryssens}
\email{wryssens@ulb.be}
\affiliation{Institut d’Astronomie et d’Astrophysique, Université Libre de 
            Bruxelles, Campus de la Plaine CP 226, 1050 Brussels, Belgium}


\author{M. Bender}
\affiliation{IPNL, Universit\'e de Lyon, Universit\'e Lyon 1, CNRS/IN2P3, 
                                                   F-69622 Villeurbanne, France}
\author{P.-H. Heenen}
\affiliation{PNTPM, CP229, Universit\'e Libre de Bruxelles, 
                                                      B-1050 Bruxelles, Belgium}
    
\date{April 2022}
\title{Brewing a cup of MOCCa: Hephaestos}


\maketitle

\setlength{\parindent}{0pt}

%%%%%%%%%%%%%%%%%%%%%%%%%%%%%%%%%%%%%%%%%%%%%%%%%%%%%%%%%%%%%%%%%%%%%%%%%%%%%%%%
\section{Directory and file structure of the entire package}


\begin{Verbatim}
 Core functionality
 -------------------
 Makefile          : contains all instructions for building working Tantalus executables.
 Hephaestos.py     : main driver for Hephaestos
 configs/          : contains predefined configuration files
 functionals/      : contains predefined functional files
 parameterizations/: contains parameterization files
 exec/             : contains executables
 mod/, obj/        : compilation directories for the compiler

 src_heph/         : Hephaestos source code
 src_orig/         : Tantalus source code, i.e. FORTRAN templates
 src/              : "completed" Tantalus source code, output of Hephaestos

 Optional/useful stuff
 -----------------------
 build_all.sh      : script to compile ALL possible configurations of Tantalus
 compilation_logs/ : compilation logs, associated with build_all.sh
 gen_XYZ_version.sh: scripts to generate public versions of the code, for 
                     distribution to various groups. 
 tantalus_examples : many examples for bugtesting and illustration
 manual            : contains Tantalus and Hephaestos manuals
 auxiliary/        : auxiliary scripts for analysing results
 
 Not-so-useful stuff
 ----------------------   
 tex/              : Folders for future interfacing with LateX documents. 
 tex_templates/    : Folders for future interfacing with LateX documents.

\end{Verbatim}

\clearpage
%%%%%%%%%%%%%%%%%%%%%%%%%%%%%%%%%%%%%%%%%%%%%%%%%%%%%%%%%%%%%%%%%%%%%%%%%%%%%%%%
\section{Running Hephaestos / Practically building Tantalus}

Hephaestos is a Python-3 code\footnote{Running with Python2.7 or before will 
likely fail.} that takes Tantalus source code templates and, depending on 
user choices, builds compileable FORTRAN code. Running Hephaestos is as simple 
as 

\begin{Verbatim}
python Hephaestos.py config
\end{Verbatim}

where \textrm{config} is the name of a configuration file from the \textrm{configs/} folder
(witthout the trailing '.py'). If everything went smoothly, you should now 
have a set of compileable Tantalus source code in the \textrm{src/} folder
that reflects the choices made in the configuration file.  

In day-to-day practice, I use the Makefile instead of running Hephaestos
manually. By typing 
%
\begin{Verbatim}
make CONFIG=config
\end{Verbatim}
%
which (i) runs Hephaestos and (ii) starts compiling immediately. Other optional
arguments for the make command are
%
\begin{Verbatim}
CXX = compiler    => ifort/gfortran work by default)
PRE = ""          => If this is given, Hephaestos is NOT run before starting the 
                     compilation process. 
\end{Verbatim}
%
At the end of a successful compilation process, you will find a new executable
in the \textrm{exec/} folder.


%%%%%%%%%%%%%%%%%%%%%%%%%%%%%%%%%%%%%%%%%%%%%%%%%%%%%%%%%%%%%%%%%%%%%%%%%%%%%%%%
\section{Configuration files}

The configuration files contain the settings with respect to symmetry choices
and advanced options regarding EDFs. The default configuration file for 
example contains the following instructions (I've removed non-relevant 
comments for brevity):  

\begin{Verbatim}
FUNC_FILE = 'NLO.func'           ! functional file to use for compilation

SYMSTRING = 'Rz,T,P,STy'         ! Symmetries to conserve in the calculation 
REDUCE    = [1,1,1]              ! Reduce (=1) or not (=0) every Cartesian axis
                                 ! [1,1,1] : calculations in 1/8th of the box
                                 ! [1,1,0] : use complete z-axis, EV4-style
                                 ! Note: the code will ensure coherence between
                                 ! SYMSTRING and this input.

INSYM     = 'Rz,T,P,STy'         ! Same options as before: these now refer to 
INREDUCE  = [1,1,1]              ! what type of .wf files this executable will
                                 ! be able to read. 
                                 
QUANT_AXIS='Z'                   ! Orientation options for the definition of
SECOND_AXIS=1                    ! multipole moments.

PH_PP_DECOUPL = True             ! If True: all feedback through density-dependencies
                                 ! of the pairing interaction will be ignored. 
                                 ! For example: a standard ULB-delta-interaction
                                 ! depends on D_I_I, but its contribution to 
                                 ! F_I_I is ignored. 
\end{Verbatim}


\clearpage
%%%%%%%%%%%%%%%%%%%%%%%%%%%%%%%%%%%%%%%%%%%%%%%%%%%%%%%%%%%%%%%%%%%%%%%%%%%%%%%%
\section{Functional files}

Functional files are the true input to Hephaestos; they contain the specification of 
the EDF. They consist of three parts, that can be described schematically as
%
\begin{itemize}
\item Description part: lines starting with '\#'. These serve to describe the 
      functional and are copied into Hephaestos output. They have no effect
      on the FORTRAN code generated. Ends with the last line starting with 
      '\#'.
\item Parameters part : a list of parameters of the functional. These are the 
      quantities that Hephaestos will assume will be read from a parameterization
      file. This is the place to put the $t_i, x_i$ and more parameters.
      Ends when the user includes a line '!TERMS'.
\item Terms part: dedicated description of terms, will be explained below.
\end{itemize}
%
In between, it is possible to write human-readable comments starting with '!'. 
These will be ignored completely by both Hephaestos and Tantalus. The only 
exception is the specific '!TERMS', which is used to indicate the transition 
between the second and third part of a .func file. 

A functional file in the end looks like this, the (abbreviated) 
content of the default functional file.

\begin{Verbatim}[fontsize=\footnotesize]
# Standard NLO Skyrme functional                                                     | Description
# -> with time-odd terms                                                             | Description
# -> without tensor terms                                                            | Description
# -> Jmunu and sDs terms optional                                                    | Description
! Particle-hole parameters                                                           | Comment
t0 ;  x0 ; t1 ;  x1 ; t2 ; x2 ; t3; x3; t3b; x3b;  wso ; wsoq ; sigma ; sigma_b      | Parameters
! Pairing parameters                                                                 | Comment
Vn ; Vp ; alpha ; sigmap  ; CT0 ; CT1 ; CSDS0 ; CSDS1                                | Parameters
!TERMS                                                                               | TRANSITION
!------------------------------------------------------------------------------------| Comment
! Central, LO terms                                                                  | Comment
! Time-even                                                                          | Comment
! Term                         ! Coupling constant            ! DD_pow! isospin      | Comment
E_D_I_I_D_I_I                  ; Cc(t0,x0,+1,0,0)             ;   1    ; 0 ; 0       | Term
E_D_I_I_D_I_I                  ; Cc(t0,x0,+1,0,1)             ;   1    ; 1 ; 1       | Term
[.......]
E_D_I_I_CP_I_I_CP_I_I          ; -Vn/4*alpha/(0.16**sigmap)   ; sigmap; 0 ; n ; n    | Term
E_D_I_I_CP_I_I_CP_I_I          ; -Vp/4*alpha/(0.16**sigmap)   ; sigmap; 0 ; p ; p    | Term
\end{Verbatim}

To include a term in the functional, one needs to specify
%
\begin{itemize}
\item Its structure.
\item Its coupling constant
\item Whether or not it is density dependent 
\item The isospin components of all densities involved. 
\end{itemize}
%
The simplest term
%
\begin{align}
C^{\rho \rho}_0 \rho_0(\vec{r}) \rho_0(\vec{r})
\end{align}
%
is encoded as
%
\begin{Verbatim}
E_D_I_I_D_I_I                  ; Cc(t0,x0,+1,0,0)                ;   1    ; 0 ; 0      
\end{Verbatim}
%
where
\begin{itemize}
\item Different columns are separated by ';'
\item The first column tells the code about $\rho \rho = D^{I,I} D^{I,I}$. 
      Note the underscores separating (i) the different operators 'I' in the densities 
      and (ii) the different densities and (iii) the initial 'E'. These are all
      mandatory.
\item The expression \textrm{ Cc(t0,x0,+1,0,0) } is the coupling constant. This
      can be any string, but Hephaestos will simply paste this into the relevant
      parts of Tantalus. Hence, this string needs to be valid FORTRAN code; 
      you can include function calls (as I do here), but these need to be 
      defined in the FORTRAN code. You can freely use any parameters included
      in the first part of the func file. Note that the convention of the 
      code for coupling constants is the most logical one: i.e. \textsf{all possible
      signs and factors} should be included in this column. I.e. the code will
      program in this case simply $E = Cc D^{I,I} D^{I,I}$.
\item The '1' indicates the power of the density dependence of the first 
      density in the term. In case this is '1', the code will ignore the 
      possibilities of this term being density-dependent.
\item The final columns indicate the isospin components: 0 or 1 for particle-hole
      terms and 'p' or 'n' for pairing densities. Note that the code expects one
      index for every density; tri/quadrilinear terms need three/four indices.
\end{itemize}

In general, the code can deal with terms $T$ of the form
%
\begin{align}
T = (X^{\hat{L}_1, \hat{R}_1}_{t_1})^{\alpha}  \prod_{i=2}^{i_{\rm max}} X^{\hat{L}_i, \hat{R}_i}_{t_i}  \, .
\end{align}
%
where $i_{\rm max} = 2,3,4$ depending if we are dealing with a 
bi-, tri- or quadrilinear term, $X$ can stand for $D, C, \tilde{D}$ or $\tilde{C}$ and
the isospin indices $t_i$ can all be independent . Hephaestos is powerful, 
but is limited by the following for the moment:
%
\begin{itemize}
\item Only the first density in the product can be raised to a power.
\item Only one vector-product $\vec{x} \times \vec{y}$ between indices 
      can be present in a single density.
\item All spin-indices/operators in all densities should be part of the 
      right-operators, i.e. in $\hat{R}_{i}$, never in the $\hat{L}_i$.
\end{itemize}
% 

Furthermore, I want to emphasize that 
\textbf{\textcolor{red}{Hephaestos is not so clever as he seems and checks
for consistency of your input only in very limited ways}}. It is \textbf{EXTREMELY
EASY} to give Hephaestos wrong input that either (i) crashes Hephaestos with 
an unintelligible error warning (ii) exits gracefully, but produces uncompileable
FORTRAN code or, even worse, (iii) generates a correct Tantalus executable that 
is subtly wrong in an unforeseen way. 

\subsection{The structure of densities}

I'll use examples to explain the writing of any given density for Hephaestos:
%
\begin{align}
D^{I,I}             &\rightarrow \textsf{D\_I\_I} \\
D^{(\nabla,\nabla)} &\rightarrow \textsf{D\_Nm\_Nm} \\
C^{I,\nabla\sigma}  &\rightarrow \textsf{C\_I\_NS} \\
\tilde{D}^{I,I}     &\rightarrow \textsf{DP\_I\_I} \, .
\end{align}
%
In short: 
\begin{itemize}
\item Replace $\nabla$ by \textsf{'N'} and $\sigma$ by \textsf{'S'}.
\item Indicate contraction by repeated summation indices like \textsf{'m'}. Other
      valid options are currently \textsf{'k', 'q', 'o', 'l'}.
\item Indicate pairing densities by adding a \textsf{'P'} after the \textsf{'D'}
      or \textsf{'C'}.     
\end{itemize}
%
To indicate external derivatives, use prefixes \textsf{Der\_} and \textsf{Lap\_}, 
such as for example
\begin{align}
\Delta D^{1,1}                  &\rightarrow \textsf{Lap\_D\_I\_I} \\
\sum_{\mu \nu \kappa} \nabla_{\mu} \epsilon_{\mu\nu\kappa} C_{\nu \kappa}^{I, \nabla \sigma}   &\rightarrow \textsf{Derm\_C\_I\_NxmSxm}\, .
\end{align}
%
This last example is perfect to demonstrate the useage of vector products: indicate
the two last indices of the Levi-civita symbol by \textsf{x}, then link these two
with the third summation index by repeated small summation letters (here \textsf{m}).

\subsection{The structure of terms}

Terms in Hephaestos are just combinations of densities linked by underscores and 
preceded by \textsf{E\_}. Of course, this means
that contractions and/or vector products can involve multiple densities. The 
code has no problems with 
%
\begin{Verbatim}
E_D_I_I_Derm_C_I_NxmSxm or E_C_Nm_NkSo_C_Nm_NkSo 
\end{Verbatim}
%
which are (i) the time-even spin-orbit term of NLO functionals and (ii) the 
$\left(\Im T_{\mu \nu \kappa}\right)^2$ terms of the Becker N2LO formulation. 


\subsection{Some remarks}

\begin{itemize}
\item The code figures out what densities and contractions to calculate from your input.
\item It time-reversal is broken is however, the cranking parts of the code require that you put at least one term with $C^{I, \nabla}$ and one with $D^{I, \sigma}$. Their coupling constants can be zero.
\item You are yourself in charge of making sure your terms are time-reversal invariant.
\item If time-reversal is conserved, the code is smart enough to drop all time-odd terms it finds.
\end{itemize}

\end{document}
